% One-time imported draft. Editable native source; no Markdown regeneration.
\section{A service request has a selector and
  arguments}\label{001-console-calculation-p05-a-service-request-has-a-selector-and-arguments}

The console interface uses register values to describe a request.
\texttt{\$v0} holds the service number. For service 4, \texttt{\$a0}
holds the address of a zero-terminated string. These two values form the
request that the next \texttt{syscall} will make.

Setting the selector prepares the request. It does not perform the
output. That separation explains a common partial-program observation:
the address and service number can both look correct in the debugger
while the console remains empty. The program must reach the request
instruction.

This is the educational MARS interface. A native executable on an
operating system has its own conventions. The syscall numbers used here
are part of the MARS environment and should be read with its service
contracts.
